\section{Conclusion and Future Work}\label{planta:sec:conclusion}

This paper presented \acronym, an ontology-mediated approach to task and architecture co-adaptation that integrates architectural knowledge represented in OWL with task planning problems defined in PDDL\@. By automatically combining these models, \acronym enables joint planning of task and reconfiguration actions using satisficing planning, with replanning when changes during execution invalidate the current plan. An open-source ROS~2 implementation of the approach is provided, including extensions to \plansys to support the PDDL features required for planning and plan execution.

The evaluation in the SUAVE underwater inspection scenario supports the feasibility of runtime co-adaptation, although it also reveals mission-performance limitations relative to other managing subsystems. In the ground navigation scenario, \acronym required less planning time and peak memory than the replicated final step of the approach by Cámara et al.~\cite{CAMARASoftwareArchitectureTask2020}, while producing plans with similar path length and estimated remaining battery. However, these approaches pursue different planning objectives, with \acronym computing satisficing plans and Cámara et al.\ generating optimal policies under probabilistic uncertainty. The scalability experiments further showed that planning time remains manageable as map size, architectural configurations, or quality requirements increase individually, but grows substantially with additional mission-level actions under optimal planning and when multiple problem dimensions are combined. In these more complex scenarios, satisficing planning remained tractable where optimal planning became impractical, although its planning time in the combined scenario remained high for applications requiring frequent replanning. These results indicate that, in the evaluated scenarios, satisficing planning with runtime replanning is a practical alternative to optimal policy generation for task and architecture co-adaptation.

Future work includes addressing the limitations discussed in \Cref{planta:sec:limitations}: selecting architectural alternatives with respect to mission-level objectives rather than per function, modeling reconfiguration costs through action costs or durative actions, and supporting numeric quality attributes without symbolic encoding. Further directions include applying \acronym to diverse real-world robotic systems to assess its reusability and further identify its limitations. Another promising direction is exploring the integration of the proposed planning approach into existing reference architectures for robot self-adaptation, such as ROSA~\cite{rosa}.

%%% Local Variables:
%%% mode: latex
%%% TeX-master: "main"
%%% End:
